\documentclass{article}
\textheight 23.5cm \textwidth 15.8cm
%\leftskip -1cm
\topmargin -1.5cm \oddsidemargin 0.3cm \evensidemargin -0.3cm
%\documentclass[final]{siamltex}

\usepackage{graphicx}
\usepackage{epsfig}
\usepackage{amsmath}
\usepackage{amsthm}
\usepackage{amssymb}
\usepackage{mathrsfs}
\usepackage{float}
\usepackage{multirow}
\usepackage{verbatim}
\usepackage{fancyhdr}
\usepackage{subfigure}
\usepackage{color}
\usepackage{mathtools}
\usepackage{mathrsfs}
%\usepackage{natbib}
\usepackage{sectsty}
%\usepackage[title]{appendix}
\usepackage{threeparttable}
\usepackage{dcolumn}
\usepackage{booktabs}
\usepackage{indentfirst}
\usepackage{setspace}
\usepackage{bm}
\usepackage{enumerate}
\usepackage{geometry}

\title{Problem 1.7: Fibonacci Search and Quadratic Interpolation}
\author{Tianyang Li}

\begin{document}
\maketitle

\section{Introduction}

This problem asks for the minimizer of the one-dimensional function
\[
    f(x) = x^4 + 2x + 4
\]
on the interval $[-1,0]$ with accuracy requirement
\[
    \epsilon = 0.01.
\]
Two numerical methods are used: Fibonacci Search and Quadratic Interpolation. The purpose is to compare the approximate minimizer obtained by these two methods and verify that both methods can achieve the required two-decimal-place accuracy.

\section{Method}

The objective function is
\[
    f(x) = x^4 + 2x + 4.
\]
Its derivative is
\[
    f'(x) = 4x^3 + 2.
\]
Setting $f'(x)=0$ gives the exact stationary point
\[
    x^* = -\sqrt[3]{\frac{1}{2}} \approx -0.793701,
\]
which lies in the interval $[-1,0]$. This value is used as a reference for the numerical results.

For Fibonacci Search, the interval is repeatedly reduced according to Fibonacci ratios. Since
\[
    \frac{b-a}{\epsilon} = \frac{1}{0.01} = 100,
\]
the first Fibonacci number not smaller than $100$ is
\[
    F_{11} = 144.
\]
Therefore the Rust implementation uses the corresponding Fibonacci sequence to determine the internal test points and finally returns the midpoint of the last interval.

For Quadratic Interpolation, three points $a$, $b$, and $c$ are maintained, and a parabola passing through $(a,f(a))$, $(b,f(b))$, and $(c,f(c))$ is constructed. Its vertex gives the next approximation
\[
    x_{new} = b - \frac{1}{2}
    \frac{(b-a)^2(f(b)-f(c)) - (b-c)^2(f(b)-f(a))}{(b-a)(f(b)-f(c)) - (b-c)(f(b)-f(a))}.
\]
The iteration stops when two consecutive approximations satisfy
\[
    |x_{new} - x_{old}| < \epsilon.
\]

\section{Result}

Running the Rust program gives the following approximations:
\[
    x_{\text{fib}} \approx -0.7917,
    \qquad
    x_{\text{quad}} \approx -0.7929.
\]
The corresponding function values are
\[
    f(x_{\text{fib}}) \approx 2.809465,
    \qquad
    f(x_{\text{quad}}) \approx 2.809452.
\]

For Fibonacci Search, the final interval is approximately
\[
    [a,b] \approx [-0.798611, -0.784722],
\]
so its midpoint gives the reported minimizer
\[
    x_{\text{fib}} = \frac{a+b}{2} \approx -0.7917.
\]
The final interval length is
\[
    b-a \approx 0.013889,
\]
which is consistent with the required accuracy.

For Quadratic Interpolation, the algorithm converges after $4$ iterations and produces
\[
    x_{\text{quad}} \approx -0.7929.
\]

To compare the two methods with the exact solution, Table~\ref{tab:compare} is given below.

\begin{table}[H]
    \centering
    \begin{tabular}{lccc}
        \toprule
        Method                  & Approximate minimizer & Function value & Absolute error \\
        \midrule
        Exact solution          & -0.793701             & 2.809449       & 0              \\
        Fibonacci Search        & -0.7917               & 2.809465       & 0.002034       \\
        Quadratic Interpolation & -0.7929               & 2.809452       & 0.000823       \\
        \bottomrule
    \end{tabular}
    \caption{Comparison of the exact solution and the two numerical methods.}
    \label{tab:compare}
\end{table}

Both approximations are correct to two decimal places:
\[
    x_{\text{fib}} \approx -0.79,
    \qquad
    x_{\text{quad}} \approx -0.79.
\]

\section{Conclusion}

For the function
\[
    f(x) = x^4 + 2x + 4
\]
on the interval $[-1,0]$, both Fibonacci Search and Quadratic Interpolation successfully locate the minimizer with accuracy better than $\epsilon=0.01$. The Fibonacci method gives
\[
    x^* \approx -0.7917,
\]
while Quadratic Interpolation gives
\[
    x^* \approx -0.7929.
\]
The exact minimizer is
\[
    x^* = -\sqrt[3]{\frac{1}{2}} \approx -0.7937.
\]
Comparing the errors shows that the Quadratic Interpolation result is slightly closer to the exact solution, but both methods meet the requirement of two-decimal-place accuracy.

\end{document}